\chapterrule
\section{从可调函数到损失：怎样比较预测？}
\label{intro:00-functions}
\label{intro:16-loss-gradient}

序章里，“下雨”“出门”“一把”让“伞”成为一个合理的后续候选。
换一种表达或上下文，合适的后续内容也可能变化。
想靠人逐条写出各种表达与语境下的预测规则，很困难。
另一个办法是选\textbf{带参数的函数}，让训练样本参与决定参数。

为了先看清参数怎样由样本确定，我们先用一个独立的数值例子练习，
后面再把这套学习思路用于文字预测。
本章与下一章始终使用三个样本 $(-1,-1),(0,2),(1,2)$：
先认识参数，再用损失比较预测；下一章学习怎样逐步改善参数。


\subsection{函数是一台按规则输出结果的机器}
\phantomsection\label{concept:function}

把一个数 $x$ 送入函数 $f$，得到输出 $y=f(x)$。例如
\[
f(x)=2x+1,\qquad f(0)=1,\qquad f(1)=3.
\]
输入也可以是一组数字。预测后续文字时，可以先把前文表示成一组数字，
再用函数为后续候选计算分数；这些步骤将在后面的章节逐步展开。
函数给我们提供了一个可计算、可检验的预测模型。

以前我们会先知道公式，再计算结果；机器学习常常反过来：先看到一些输入与结果，再寻找一个能够解释这些例子的函数。

\subsection{三个样本，一条可调的直线}
\phantomsection\label{concept:sample}

设训练数据是
\[
(-1,-1),\quad(0,2),\quad(1,2).
\]
我们先选一个最简单的函数族：
\begin{equation}
\widehat{y}=wx+b.
\label{eq:line-first}
\end{equation}
$\widehat{y}$ 表示模型的\textbf{预测值}，而 $y$ 是实际观察到的值。$w$ 决定斜率，$b$ 决定截距。

先令 $w=1,b=0$，模型成为 $\widehat{y}=x$。
三个输入依次得到 $(-1,0,1)$，与真实值 $(-1,2,2)$ 比较，残差为 $(0,-2,-1)$。
只有第一点预测正确。保持 $w=1$，把 $b$ 从 $0$ 调到 $1$，预测变成
\[
\widehat{y}(-1)=0,\qquad \widehat{y}(0)=1,\qquad \widehat{y}(1)=2.
\]
第二、三个样本的预测更接近真实值了，第一个样本却从猜对变成了猜错。
若再把 $b$ 加到 $2$，中间的样本猜对了，两端的预测又都偏高。
参数调整会同时影响多个样本，需要比较整组预测的表现。

为什么两个参数各有用处？增大 $b$ 会把三个预测一起向上移；
增大 $w$ 则让 $x=-1$ 的预测下降、$x=1$ 的预测上升，$x=0$ 的预测保持不变。
这三个点也无法被同一条直线全部穿过：输入每次增加 $1$，
输出却先增加 $3$、再增加 $0$，而直线每次的输出增量相同。
我们要寻找一组能让整体预测更合适的参数。

\par\noindent\begin{minipage}{\linewidth}
\tagstructbegin{tag=Figure,alttext={三个不共线的观察点负一负一、零二、一二。三条候选直线分别为 y 等于 x、x 加一、x 加二；增大偏置使直线整体上移，某个点更准时另一个点可能变差，没有一条直线穿过全部三点。}}
\begin{center}
\begin{tikzpicture}
\begin{axis}[
width=.77\linewidth,height=57mm,
axis lines=middle,xmin=-1.7,xmax=1.7,ymin=-2.2,ymax=4.2,
xlabel={$x$},ylabel={$y$},grid=both,
grid style={linegray,thin},legend pos=north west,
legend style={font=\footnotesize},
xtick={-1,0,1},ytick={-1,0,1,2,3},
]
\addplot[datacolor,thick,dashed,domain=-1.4:1.35,samples=2]{x};
\addlegendentry{候选 A：$\hat{y}=x$}
\addplot[datacolor,very thick,domain=-1.4:1.35,samples=2]{x+1};
\addlegendentry{候选 B：$\hat{y}=x+1$}
\addplot[datacolor,thick,dashdotted,domain=-1.4:1.35,samples=2]{x+2};
\addlegendentry{候选 C：$\hat{y}=x+2$}
\addplot[only marks,mark=*,noisecolor,mark size=2.5pt] coordinates {(-1,-1) (0,2) (1,2)};
\addlegendentry{观察数据}
\end{axis}
\end{tikzpicture}
\captionof{figure}{同一组三点与三条候选直线：调整偏置会同时改变各个样本的预测}
\label{fig:first-line-fit}
\end{center}
\tagstructend
\end{minipage}\par
\medskip

\subsection{从可调函数到模型、训练与推理}
\phantomsection\label{concept:model}
\phantomsection\label{concept:parameter}
\phantomsection\label{concept:weight}
\phantomsection\label{concept:training}
\phantomsection\label{concept:inference}

请先分清四个词：
\begin{enumerate}
\item \textbf{模型（model）}：我们选择的函数形式，例如 $\widehat{y}=wx+b$。
\item \textbf{参数（parameters）}：模型内部可调的数，例如 $w,b$。
\item \textbf{训练（training）}：根据样本调整参数，让预测更符合任务目标。
\item \textbf{推理（inference）}：参数确定后，输入新数据，计算预测结果。
\end{enumerate}
在深度学习里，函数可以有数百万乃至更多的可训练参数；但计算过程仍是\textbf{输入、运算、输出}。

拟合需要用新数据检验：模型在已有样本上表现好，对新数据仍可能预测不准。后续的\textbf{泛化}章节专门讨论这个问题。



\subsection{正负残差相抵，就算猜对了吗？}
\phantomsection\label{concept:loss}

沿用三个样本 $(-1,-1),(0,2),(1,2)$ 与模型 $\widehat{y}=wx+b$。
图~\ref{fig:first-line-fit} 中的候选 B 取 $w=1,b=1$，预测为 $(0,1,2)$。
每个样本的\textbf{残差}是预测减真实：$\widehat{y}-y$，三项分别为
\[
0-(-1)=1,\qquad 1-2=-1,\qquad 2-2=0.
\]
\phantomsection\label{partq:learning}
把它们相加，得到 $1+(-1)+0=0$。这回答了本部分开篇的问题：
前两个点分别猜高和猜低，偏差相抵，却都没有猜对。
我们需要让每个样本的偏差都计入分数。

\subsection{平方，再平均：均方误差}

先把三个残差各自平方，再取平均。候选 B 的分数为
\[
\frac{1^2+(-1)^2+0^2}{3}=\frac23.
\]
这叫\textbf{均方误差}（Mean Squared Error，MSE），是一种\textbf{损失函数}：
用一个数衡量整组预测与真实值的差距，数值越小表示按这项标准预测得越好。
对当前三个样本，写成
\begin{equation}
\boxed{L(w,b)=\frac13\sum_{i=1}^{3}(\widehat{y}_i-y_i)^2.}
\label{eq:three-point-mse}
\end{equation}
这里 $i$ 是样本编号，$\widehat{y}_i=wx_i+b$。
平方使正负残差都成为非负数，也让较大的残差受到更大惩罚；
取平均则得到每个样本的平均平方误差。

\subsection{用同一把尺子比较三组参数}

从候选 A 的 $w=1,b=0$ 出发，三项预测为 $(-1,0,1)$，残差为 $(0,-2,-1)$，所以
\[
L(1,0)=\frac{0^2+(-2)^2+(-1)^2}{3}=\boxed{\frac53}.
\]
把三条候选直线放在同一张表中：
\begin{center}
\tagpdfsetup{table/header-rows={1}}
\begin{tabular}{lcccc}
\toprule
候选 & $(w,b)$ & 三项预测 & 三项残差 & 均方误差\\\midrule
A & $(1,0)$ & $(-1,0,1)$ & $(0,-2,-1)$ & $5/3$\\
B & $(1,1)$ & $(0,1,2)$ & $(1,-1,0)$ & $2/3$\\
C & $(1,2)$ & $(1,2,3)$ & $(2,0,1)$ & $5/3$\\
\bottomrule
\end{tabular}
\end{center}
\textbf{B 在这三个候选中损失最小。}C 把中间的样本猜对了，
两端却猜得更差，整体损失反而比 B 大。
改善某一个样本的预测，可能要付出其他样本误差增大的代价；损失把这些影响一起计入。

\subsection{有了评分标准，还需要寻找参数的方法}

三个点不共线，一条直线无法让所有残差同时归零。
在选定的直线模型与均方误差下，我们寻找的是整体误差尽量小的折中。
B 胜过了 A、C，但还有无数组 $w,b$ 尚未比较，是否存在损失更小的参数？

为了下一章计算方便，把同一个损失按样本写全：
\begin{equation}
L(w,b)=\frac{(-w+b+1)^2+(b-2)^2+(w+b-2)^2}{3}.
\label{eq:three-point-residual-loss}
\end{equation}
三个平方项按顺序对应 $x=-1,0,1$ 的样本。
给定一组参数，我们已经能计算它的分数；接下来要让当前误差帮助我们决定怎样调整参数。

\chapterchecktasks{00-functions}{理解、手算与判断}{
\label{check:16-loss-analytic}
\item\textbf{预测：}取 $y=wx+b$，固定 $x=3$、$b=1$，把 $w$ 从 $2$ 增到 $3$，输出怎样变化？先写变化量，再算新输出。\chapterinlinehint{固定输入，先算权重增量乘以输入。}
\item\textbf{找错：}同学说“推理时输入变了，训练好的 $w,b$ 也会自动再训练一遍”。请分别指出训练与推理究竟改变什么。\chapterinlinehint{分别列出训练阶段和推理阶段的可变量。}
\item\textbf{解释：}候选 B 的残差为 $(1,-1,0)$，相加为零。为什么均方误差仍大于零？\chapterinlinehint{先分别判断前两个样本是否预测正确，再平方取平均。}
\item\textbf{手算：}沿用本章三个样本，取另一组参数 $w=2,b=1$。求三项预测、残差与均方误差，并与候选 B 比较。\chapterinlinehint{按 $x=-1,0,1$ 的顺序计算，不只看猜对了几个点。}
\item\textbf{找错：}有人说“B 的损失比 A、C 都小，所以 B 一定是所有直线中的最优解”。这个判断缺少什么依据？\chapterinlinehint{区分三个候选中的最小值与全部参数中的最小值。}
}

\par\smallskip
损失给整组预测一个可比较的分数，均方误差把各个样本的残差平方后取平均。
下一章从当前参数出发，用损失的变化方向寻找更好的参数。
